% LNCS: tóm tắt 150--250 từ (tính trên bản tiếng Anh).
\begin{abstract}
Tác tử LLM theo mô hình ReAct đọc nội dung từ các công cụ bên ngoài. Chỉ thị được cài vào nội dung đó có thể khiến tác
tử thực hiện hành động trái phép hoặc làm lộ dữ liệu. Đây là tấn công chèn lệnh gián tiếp (IPI). Các benchmark an
ninh cho tác tử hiện nay hầu như chỉ có tiếng Anh. Bài báo đề xuất \ours, một cơ chế phòng thủ nhiều lớp đặt việc
cưỡng chế tại Tool Broker, bên ngoài mô hình. Các luật của \ours{} chỉ cho phép gửi tới đích do người dùng viết, chỉ
cho phép giá trị mật rời tác tử khi người dùng đã viết giá trị đó, che giá trị mật trong câu trả lời cuối và chặn
lời gọi tới đích được nêu trong nội dung có dấu hiệu chèn lệnh. Lời nhắc trong prompt và phản hồi từ chối bổ trợ cho
các luật. \ours{} không cần tinh chỉnh mô hình và không gọi thêm LLM. Chúng tôi xây dựng VN-AgentInject v3, một
benchmark tiếng Việt gồm sáu cơ chế tấn công, có nhánh lành tính và năm biến thể ngôn ngữ, và đánh giá theo giao thức
đăng ký trước với hai tập kiểm thử niêm phong. Trên ba mô hình nền mở trọng số, phiên bản thứ nhất (\cRulesS) giảm tỉ
lệ tấn công thành công từ 32,9\% xuống 1,2\% trên tập niêm phong thứ nhất. Phiên bản thứ hai (\bgTwo) giảm tỉ lệ này
từ 32,4\% xuống 3,6\% trên tập niêm phong thứ hai (Holm $p = 5{,}0\cdot10^{-57}$) và giữ 93,7\% tiện ích lành tính.
Bộ phát hiện huấn luyện trên tiếng Anh không có tác dụng trên tiếng Việt.

\keywords{Chèn lệnh gián tiếp \and Tác tử LLM \and An ninh tác tử \and Kiểm soát luồng thông tin \and Tiếng Việt}
\end{abstract}
